\subsection{Divided powers}

Let \(x \in M\) and \(k \in N\) . It is clear that \(x_1 \otimes x_2 \otimes \ldots \otimes x_k\) , where

\[
x _ {1} = x _ {2} = \dots = x _ {k} = x,
\]

is an element of \(\mathbf{TS}^k (\mathbf{M})\)

DEFINITION 2. --- If \(x \in \mathbf{M}\) , the element \(x \otimes x \otimes \ldots \otimes x\) of \(\mathbf{T}\mathbf{S}^k(\mathbf{M})\) is denoted by \(\gamma_k(x)\) .

PROPOSITION 3. --- (i) If \(x \in \mathbf{M}\) , the \(p\) th power of \(x\) , calculated in \(\mathbf{T}\mathbf{S}(\mathbf{M})\) , is equal to \(p! \gamma_p(x)\) .

\begin{enumerate}
\def\labelenumi{(\roman{enumi})}
\setcounter{enumi}{1}
\tightlist
\item
  Let \(x_{1},\ldots ,x_{n}\in \mathbf{M};\) then
\end{enumerate}

\[
\gamma_ {p} \left(x _ {1} + x _ {2} + \dots + x _ {n}\right) = \sum_ {p _ {1} + p _ {2} + \dots + p _ {n} = p} \gamma_ {p _ {1}} \left(x _ {1}\right) \gamma_ {p _ {2}} \left(x _ {2}\right) \dots \gamma_ {p _ {n}} \left(x _ {n}\right).
\]

\begin{enumerate}
\def\labelenumi{(\roman{enumi})}
\setcounter{enumi}{2}
\tightlist
\item
  Let \(x_1, \ldots, x_n \in \mathbf{M}\) , let \(p_1, \ldots, p_n\) be integers \(\geqslant 0\) and \(p = p_1 + \cdots + p_n\) . Let \(E\) be the set of mappings \(\varphi\) of \(\{1, \ldots, p\}\) into \(\{1, \ldots, n\}\) such that
\end{enumerate}

\[
\operatorname{Card} \varphi^ {- 1} (1) = p _ {1}, \dots , \operatorname{Card} \varphi^ {- 1} (n) = p _ {n}.
\]

Then

\[
\gamma_ {p _ {1}} (x _ {1}) \gamma_ {p _ {2}} (x _ {2}) \dots \gamma_ {p _ {n}} (x _ {n}) = \sum_ {\varphi \in E} x _ {\varphi (1)} \otimes x _ {\varphi (2)} \otimes \dots \otimes x _ {\varphi (p)}.
\]

\begin{enumerate}
\def\labelenumi{(\roman{enumi})}
\setcounter{enumi}{3}
\tightlist
\item
  Let \(x \in \mathbf{M}\) and let \(q, r\) be integers \(\geqslant 0\) . Then
\end{enumerate}

\[
\gamma_ {q} (x) \gamma_ {r} (x) = \frac {(q + r) !}{q ! r !} \gamma_ {q + r} (x).
\]

\begin{enumerate}
\def\labelenumi{(\alph{enumi})}
\setcounter{enumi}{21}
\tightlist
\item
  Let \(x_{1}, \ldots, x_{n} \in \mathbf{M}\) , and for \(\mathrm{H} \subset \{1, \ldots, n\}\) put \(x_{\mathrm{H}} = \sum_{i \in \mathrm{H}} x_{i}\) , then
\end{enumerate}

\[
(- 1) ^ {n} x _ {1} x _ {2} \dots x _ {n} = \sum_ {\mathrm{H} \subset \{1, \dots , n \}} (- 1) ^ {\text { Card   H }} \gamma_ {n} (x _ {\mathrm{H}}).
\]

The assertion (i) follows at once from Prop. 2 (ii).

Let us prove (ii) ; by an induction on \(n\) we see that it is enough to consider the case \(n = 2\) . Then we have

\[
\begin{array}{l} \gamma_ {p} (x _ {1} + x _ {2}) = (x _ {1} + x _ {2}) \otimes (x _ {1} + x _ {2}) \otimes \dots \otimes (x _ {1} + x _ {2}) \quad (p \text {factors}) \\ = \sum_ {p _ {1} + p _ {2} = p} \sum_ {\sigma \in \mathfrak {S} _ {p _ {1}, p _ {2}}} \sigma (x _ {1} \otimes x _ {1} \otimes \dots \otimes x _ {1} \otimes x _ {2} \otimes x _ {2} \otimes \dots \otimes x _ {2}) \\ = \sum_ {p _ {1} + p _ {2} = p} \sum_ {\sigma \in \mathfrak {S} _ {p _ {1}, p _ {2}}} \sigma (\gamma_ {p _ {1}} (x _ {1}) \otimes \gamma_ {p _ {2}} (x _ {2})) \\ = \sum_ {p _ {1} + p _ {2} = p} \gamma_ {p _ {1}} (x _ {1}) \gamma_ {p _ {2}} (x _ {2}). \end{array}
\]

To prove (iii), let \(\mathfrak{S}_{p_{1},\ldots,p_{n}}\) be the set of permutations of \((1,p_{1}+\cdots+p_{n})\) whose restrictions to the intervals

\[
\left(1, p _ {1}\right), \left(p _ {1} + 1, p _ {1} + p _ {2}\right), \dots , \left(p _ {1} + \dots + p _ {n - 1} + 1, p _ {1} + \dots + p _ {n}\right)
\]

are increasing. By I, p.~60, example 2 and Prop. 2, (ii) we have

\[
\gamma_ {p _ {1}} \left(x _ {1}\right) \gamma_ {p _ {2}} \left(x _ {2}\right) \dots \gamma_ {p _ {n}} \left(x _ {n}\right) = \sum_ {\rho \in \mathfrak {S} _ {p _ {1}, p _ {2}, \dots , p _ {n}}} \rho \left(x _ {1} \otimes x _ {1} \otimes \dots \otimes x _ {1} \otimes x _ {2} \otimes x _ {2} \otimes \dots \otimes x _ {2} \otimes \dots \otimes x _ {n} \otimes x _ {n} \otimes x _ {n} \otimes \dots \otimes x _ {n}\right)
\]

(with \(p_i\) factors \(x_i\) ) and this sum is equal to

\[
\sum_ {\varphi \in E} x _ {\varphi (1)} \otimes x _ {\varphi (2)} \otimes \dots \otimes x _ {\varphi (p)}.
\]

In (iii) let us put \(n = 2\) , \(x_{1} = x_{2} = x\) , \(p_{1} = q\) and \(p_{2} = r\) , then we obtain (iv) (I, loc. cit.).

Finally (v) follows from Prop. 2, (ii) and Prop. 2 of I, p.~100, applied to the elements \(x_{i}\) of the ring \(\mathbf{T}(\mathbf{M})\) .

Remarks. --- 1) Let \((x_{i})_{i\in I}\) be a family of elements of M. For each \(\nu \in \mathbf{N}^{(I)}\) put

\[
x _ {\nu} = \prod_ {i \in I} \gamma_ {\nu_ {i}} (x _ {i}).
\]

If \((\lambda_{i})\in \mathbf{A}^{(\mathrm{I})}\) and \(p\in \mathbf{N}\) , then we have by Prop. 3 (ii),

\[
\gamma_ {p} \left(\sum_ {i \in I} \lambda_ {i} x _ {i}\right) = \sum_ {\nu \in \mathbf {N} ^ {(I)}, | \nu | = p} \lambda^ {\nu} x _ {\nu}.\tag{6}
\]

\begin{enumerate}
\def\labelenumi{\arabic{enumi})}
\setcounter{enumi}{1}
\tightlist
\item
  Let \(\mathcal{M}\) be the set of mappings of \([1, p]\) into I. We define a mapping \(\rho \mapsto \rho^{*}\) of \(\mathcal{M}\) into \(\mathbf{N}^{(I)}\) by putting
\end{enumerate}

\[
\rho^ {*} (i) = \operatorname{Card} \rho^ {- 1} (i).
\]

For two elements \(\rho_1, \rho_2\) of \(\mathcal{M}\) to satisfy \(\rho_1^* = \rho_2^*\) it is necessary and sufficient that there should exist \(\sigma \in \mathfrak{S}_p\) such that \(\rho_2 = \rho_1 \circ \sigma\) (I, p.~95). By Prop. 3 (iii) we have, for \(|\nu| = p\) ,

\[
x _ {\nu} = \sum_ {\rho \in \mathcal {M}, \rho^ {*} = \nu} x _ {\rho (1)} \otimes x _ {\rho (2)} \otimes \dots \otimes x _ {\rho (p)}.\tag{7}
\]

